\documentclass[11pt,a4paper]{article}

% Packages
\usepackage[utf8]{inputenc}
\usepackage[margin=1in]{geometry}
\usepackage{hyperref}
\usepackage{xcolor}
\usepackage{listings}
\usepackage{booktabs}
\usepackage{amsmath,amssymb}
\usepackage{fancyhdr}
\usepackage{enumitem}

% Colors for code snippets
\definecolor{codegreen}{rgb}{0,0.6,0}
\definecolor{codegray}{rgb}{0.5,0.5,0.5}
\definecolor{codepurple}{rgb}{0.58,0.82,0.82}
\definecolor{backcolour}{rgb}{0.96,0.96,0.96}
\definecolor{navy}{rgb}{0.0,0.2,0.6}

% Listing setup
\lstdefinestyle{mystyle}{
    backgroundcolor=\color{backcolour},   
    commentstyle=\color{codegreen},
    keywordstyle=\color{navy}\bfseries,
    numberstyle=\tiny\color{codegray},
    stringstyle=\color{codepurple},
    basicstyle=\ttfamily\footnotesize,
    breakatwhitespace=false,         
    breaklines=true,                 
    captionpos=b,                    
    keepspaces=true,                 
    numbers=left,                    
    numbersep=5pt,                  
    showspaces=false,                
    showstringspaces=false,
    showtabs=false,                  
    tabsize=2
}
\lstset{style=mystyle}

% Page Header/Footer
\pagestyle{fancy}
\fancyhf{}
\rhead{\textbf{Crazyflie 2.1 ROS 2 Integration}}
\lhead{Technical Documentation}
\rfoot{Page \thepage}

\title{\textbf{\Huge Crazyflie 2.1 ROS 2 Humble Integration \& Modular Control Nodes Document}\\[0.5em]
\Large Complete Setup, Node Architecture, and Battery Monitoring}
\author{\textbf{Karan Kumar Anand} \\ Department of Computer Science \& Engineering \\ IIIT Delhi}
\date{July 27, 2026}

\begin{document}

\maketitle

\begin{abstract}
This document details the software architecture, modular node development, system configuration, and battery diagnostic tools implemented for the Crazyflie 2.1 micro-quadcopter within a ROS 2 Humble environment. It covers the \texttt{my\_crazyflie} ROS 2 Python package containing specialized clients for arming, takeoff, landing, emergency stop, relative movement, and velocity control, as well as a standalone \texttt{cflib}-based battery monitoring utility.
\end{abstract}

\tableofcontents
\newpage

\section{Introduction \& Workspace Overview}
The Crazyflie 2.1 workspace is organized under \texttt{\~/cf\_ws} operating on ROS 2 Humble Linux. The hardware platform used is a stock Crazyflie 2.1 quadcopter configured with zero expansion decks (0-deck setup). Due to the lack of optical flow or positioning decks, position estimation relies on open-loop or external positioning methods.

\subsection{Environment Setup}
Before running ROS 2 nodes or launching the Crazyflie server, the environment must be sourced:
\begin{lstlisting}[language=bash]
cd ~/cf_ws
source /opt/ros/humble/setup.bash
source install/setup.bash
\end{lstlisting}

---

\section{ROS 2 Package Structure: \texttt{my\_crazyflie}}
A custom ROS 2 Python package named \texttt{my\_crazyflie} was constructed in \texttt{\~/cf\_ws/src/my\_crazyflie}.

\subsection{Package Entry Points (\texttt{setup.py})}
The console scripts registered in \texttt{setup.py} allow direct execution via \texttt{ros2 run my\_crazyflie <script>}:
\begin{lstlisting}[language=python]
entry_points={
    'console_scripts': [
        'takeoff          = my_crazyflie.takeoff:main',
        'land             = my_crazyflie.land:main',
        'arm              = my_crazyflie.arm:main',
        'emergency        = my_crazyflie.emergency:main',
        'sequence         = my_crazyflie.sequence:main',
        'go_to            = my_crazyflie.go_to:main',
        'velocity_control = my_crazyflie.velocity_control:main',
    ],
},
\end{lstlisting}

---

\section{Detailed Control Nodes Implementation}

\subsection{\texttt{arm.py} — Drone Arming Node}
In CrazySwarm2, arming is required before executing takeoff commands. This client uses \texttt{crazyflie\_interfaces/srv/Arm}.

\begin{lstlisting}[language=python]
import rclpy
from rclpy.node import Node
from crazyflie_interfaces.srv import Arm

class ArmClient(Node):
    def __init__(self):
        super().__init__("arm_client")
        self.client = self.create_client(Arm, "/cf231/arm")
        while not self.client.wait_for_service(timeout_sec=1.0):
            self.get_logger().info("Waiting for arm service...")

    def arm(self):
        request = Arm.Request()
        request.arm = True
        self.get_logger().info("Arming drone...")
        future = self.client.call_async(request)
        rclpy.spin_until_future_complete(self, future)
        self.get_logger().info("Drone ARMED!")
\end{lstlisting}

\subsection{\texttt{takeoff.py} — Autonomous Takeoff Node}
Sends a height target and duration to \texttt{/cf231/takeoff}.

\begin{lstlisting}[language=python]
import rclpy
from rclpy.node import Node
from crazyflie_interfaces.srv import Takeoff

class TakeoffClient(Node):
    def __init__(self):
        super().__init__("takeoff_client")
        self.client = self.create_client(Takeoff, "/cf231/takeoff")

    def send_request(self):
        request = Takeoff.Request()
        request.group_mask = 0
        request.height = 0.2
        request.duration.sec = 2
        future = self.client.call_async(request)
        rclpy.spin_until_future_complete(self, future)
\end{lstlisting}

\subsection{\texttt{land.py} — Controlled Landing Node}
Commands the drone to smoothly descent to $0.0\text{ m}$ height over $2\text{ seconds}$.

\subsection{\texttt{emergency.py} — Emergency Stop Node}
Calls \texttt{/cf231/emergency} (\texttt{std\_srvs/srv/Empty}) to cut motor power immediately in case of unexpected flight drift.

\subsection{\texttt{sequence.py} — Complete Automated Flight Routine}
Executes the full safe flight state machine:
$$\text{Arm} \longrightarrow \text{Takeoff (0.3m)} \longrightarrow \text{Hover (3s)} \longrightarrow \text{Land (0.0m)}$$

\subsection{\texttt{go\_to.py} \& \texttt{velocity_control.py}}
\begin{itemize}
    \item \texttt{go\_to.py}: Uses \texttt{crazyflie\_interfaces/srv/GoTo} to move to relative coordinates $(x, y, z)$.
    \item \texttt{velocity\_control.py}: Publishes velocity and attitude commands directly over \texttt{/cf231/cmd\_vel_legacy} using \texttt{geometry\_msgs/msg/Twist}.
\end{itemize}

---

\section{Crazyflie Battery Diagnostics (\texttt{battery\_check.py})}
When using zero expansion decks, power management is critical. Low battery levels result in motor spin without sufficient thrust for lift-off.

\subsection{Standalone Diagnostics Script}
A Python script using the native \texttt{cflib} library was developed to poll battery status directly via Crazyradio without needing the ROS server active.

\subsection{1S LiPo Battery Voltage Reference Table}
The Crazyflie 2.1 uses a 1S 3.7V nominal LiPo battery ($3.0\text{ V}$ cutoff to $4.2\text{ V}$ max).

\begin{table}[h!]
\centering
\caption{Crazyflie 2.1 LiPo 1S Voltage Classification}
\vspace{0.5em}
\begin{tabular}{@{}ccll@{}}
\toprule
\textbf{Voltage Range (V)} & \textbf{Approx. Charge (\%)} & \textbf{Status Indicator} & \textbf{Recommended Action} \\ \midrule
$\ge 4.10\text{ V}$ & $\sim 100\%$ & 🟢 Full & Ready for flight \\
$\ge 3.90\text{ V}$ & $\sim 75\%$ & 🟡 Good & Safe for normal testing \\
$\ge 3.75\text{ V}$ & $\sim 50\%$ & 🟠 Medium & Short flight tests only \\
$\ge 3.50\text{ V}$ & $\sim 25\%$ & 🔴 Low & Recharge soon \\
$\ge 3.20\text{ V}$ & $\sim 10\%$ & ⛔ Very Low & Land \& recharge immediately \\
$< 3.20\text{ V}$ & $< 5\%$ & 💀 Critical & Do NOT operate (Battery protection) \\ \bottomrule
\end{tabular}
\end{table}

\subsection{Executing Battery Check}
Ensure the Crazyflie ROS server is stopped, then run:
\begin{lstlisting}[language=bash]
python3 src/my_crazyflie/my_crazyflie/battery_check.py
\end{lstlisting}

---

\section{Command Summary \& Verification Procedures}

\subsection{Step 1: Start Server (Terminal 1)}
\begin{lstlisting}[language=bash]
cd ~/cf_ws
source /opt/ros/humble/setup.bash
source install/setup.bash
ros2 launch crazyflie launch.py
\end{lstlisting}

\subsection{Step 2: Build Workspace \& Run Flight Sequence (Terminal 2)}
\begin{lstlisting}[language=bash]
cd ~/cf_ws
source /opt/ros/humble/setup.bash
colcon build --symlink-install
source install/setup.bash

# Run automated flight sequence
ros2 run my_crazyflie sequence
\end{lstlisting}

\section{Conclusion}
The integration of ROS 2 Humble with the Crazyflie 2.1 micro-quadcopter under \texttt{\~/cf\_ws} has been successfully verified. With arming interfaces, automated flight sequences, and battery diagnostic scripts implemented, the system is fully operational for autonomous aerial robotics research.

\end{document}
